\documentclass[10pt,a4paper]{amsart}
\usepackage{amsmath,amssymb,amsthm,graphicx}
\usepackage{verbatim}
\newenvironment{reference}{\comment}{\endcomment}
%\newenvironment{reference}{}{}
\newenvironment{slogan}{\comment}{\endcomment}
\newenvironment{history}{\comment}{\endcomment}

% For commutative diagrams we use Xy-pic
\usepackage[all]{xy}

% We use 2cell for 2-commutative diagrams.
\xyoption{2cell}
% Theorem environments.
%
\theoremstyle{plain}
\newtheorem{theorem}[subsection]{Theorem}
\newtheorem{proposition}[subsection]{Proposition}
\newtheorem{lemma}[subsection]{Lemma}

\theoremstyle{definition}
\newtheorem{definition}[subsection]{Definition}
\newtheorem{example}[subsection]{Example}
\newtheorem{exercise}[subsection]{Exercise}
\newtheorem{situation}[subsection]{Situation}

\theoremstyle{remark}
\newtheorem{remark}[subsection]{Remark}
\newtheorem{remarks}[subsection]{Remarks}

\numberwithin{equation}{subsection}

% Macros
%
\def\lim{\mathop{\mathrm{lim}}\nolimits}
\def\colim{\mathop{\mathrm{colim}}\nolimits}
\def\Spec{\mathop{\mathrm{Spec}}}
\def\Hom{\mathop{\mathrm{Hom}}\nolimits}
\def\Ext{\mathop{\mathrm{Ext}}\nolimits}
\def\SheafHom{\mathop{\mathcal{H}\!\mathit{om}}\nolimits}
\def\SheafExt{\mathop{\mathcal{E}\!\mathit{xt}}\nolimits}
\def\Sch{\mathit{Sch}}
\def\Mor{\mathop{\mathrm{Mor}}\nolimits}
\def\Ob{\mathop{\mathrm{Ob}}\nolimits}
\def\Sh{\mathop{\mathit{Sh}}\nolimits}
\def\NL{\mathop{N\!L}\nolimits}
\def\CH{\mathop{\mathrm{CH}}\nolimits}
\def\proetale{{pro\text{-}\acute{e}tale}}
\def\etale{{\acute{e}tale}}
\def\QCoh{\mathit{QCoh}}
\def\Ker{\mathop{\mathrm{Ker}}}
\def\Im{\mathop{\mathrm{Im}}}
\def\Coker{\mathop{\mathrm{Coker}}}
\def\Coim{\mathop{\mathrm{Coim}}}

% Boxtimes
%
\DeclareMathSymbol{\boxtimes}{\mathbin}{AMSa}{"02}

%
% Macros for moduli stacks/spaces
%
\def\QCohstack{\mathcal{QC}\!\mathit{oh}}
\def\Cohstack{\mathcal{C}\!\mathit{oh}}
\def\Spacesstack{\mathcal{S}\!\mathit{paces}}
\def\Quotfunctor{\mathrm{Quot}}
\def\Hilbfunctor{\mathrm{Hilb}}
\def\Curvesstack{\mathcal{C}\!\mathit{urves}}
\def\Polarizedstack{\mathcal{P}\!\mathit{olarized}}
\def\Complexesstack{\mathcal{C}\!\mathit{omplexes}}
% \Pic is the operator that assigns to X its picard group, usage \Pic(X)
% \Picardstack_{X/B} denotes the Picard stack of X over B
% \Picardfunctor_{X/B} denotes the Picard functor of X over B
\def\Pic{\mathop{\mathrm{Pic}}\nolimits}
\def\Picardstack{\mathcal{P}\!\mathit{ic}}
\def\Picardfunctor{\mathrm{Pic}}
\def\Deformationcategory{\mathcal{D}\!\mathit{ef}}
\usepackage{cleveref}
\crefname{originallemma}{original lemma}{original lemmas}
\Crefname{originallemma}{Original lemma}{Original lemmas}
\makeatletter
\newlabel{algebra-lemma-locally-free-semi-local-free}{{02M9}{source}{}{sourcepointer-1}{}}
\newlabel{algebra-lemma-locally-free-semi-local-free@cref}{{[originallemma][0][]02M9}{[1][0][]source}{}{}{}}
\makeatother
\crefname{originallemma}{original lemma}{original lemmas}
\Crefname{originallemma}{Original lemma}{Original lemmas}
\makeatletter
\newlabel{divisors-lemma-effective-Cartier-makes-dimension-drop}{{056N}{source}{}{sourcepointer-2}{}}
\newlabel{divisors-lemma-effective-Cartier-makes-dimension-drop@cref}{{[originallemma][0][]056N}{[1][0][]source}{}{}{}}
\makeatother
\crefname{originalsection}{original section}{original sections}
\Crefname{originalsection}{Original section}{Original sections}
\makeatletter
\newlabel{divisors-section-effective-Cartier-invertible}{{0C4S}{source}{}{sourcepointer-3}{}}
\newlabel{divisors-section-effective-Cartier-invertible@cref}{{[originalsection][0][]0C4S}{[1][0][]source}{}{}{}}
\makeatother
\crefname{originallemma}{original lemma}{original lemmas}
\Crefname{originallemma}{Original lemma}{Original lemmas}
\makeatletter
\newlabel{morphisms-lemma-proper-ample-is-proj}{{0C6J}{source}{}{sourcepointer-4}{}}
\newlabel{morphisms-lemma-proper-ample-is-proj@cref}{{[originallemma][0][]0C6J}{[1][0][]source}{}{}{}}
\makeatother
\crefname{originallemma}{original lemma}{original lemmas}
\Crefname{originallemma}{Original lemma}{Original lemmas}
\makeatletter
\newlabel{varieties-lemma-algebraic-scheme-dim-0}{{06LH}{source}{}{sourcepointer-5}{}}
\newlabel{varieties-lemma-algebraic-scheme-dim-0@cref}{{[originallemma][0][]06LH}{[1][0][]source}{}{}{}}
\makeatother
\crefname{originallemma}{original lemma}{original lemmas}
\Crefname{originallemma}{Original lemma}{Original lemmas}
\makeatletter
\newlabel{varieties-lemma-chi-tensor-finite}{{0AYT}{source}{}{sourcepointer-6}{}}
\newlabel{varieties-lemma-chi-tensor-finite@cref}{{[originallemma][0][]0AYT}{[1][0][]source}{}{}{}}
\makeatother
\title{On a proper curve with H\textasciicircum{}0(X, O\_X) = k, a globally generated-in-principle invertible sheaf with a regular section, vanishing H\textasciicircum{}1 and ampleness has very ample sixth tensor power}
\author{The Stacks Project Authors}
\date{stacks-project:curves.tex:local-snapshot-20261012; GNU FDL 1.2 or later}
\begin{document}
\maketitle
\noindent\textit{Editorial note.} All delivered mathematical text is the author's own, copied from the pinned Stacks Project chapter file curves.tex; no statement or proof step was rewritten, reordered or removed.\par\medskip
\noindent\textit{Editorial note.} The counted proof is the author's complete argument: the Cartier zero scheme setup with the author's small-detail remark on nonemptiness, the induction giving vanishing and trivializations, the two-term multiplication map surjectivity, and the closing Veronese generation argument are all retained verbatim.\par\medskip
\noindent\textit{Editorial note.} The label lemma-criterion-very-ample names the counted statement, and the proof environment immediately following it in the source is delivered verbatim as its proof.\par\medskip
\noindent\textit{Editorial note.} The counted statement is carried by the author's own lemma environment in the chapter source.\par\medskip
\noindent\textit{Allowed standard prerequisites (editorial).}\begin{itemize}
\item Zero schemes of regular sections. Divisors, Section divisors-section-effective-Cartier-invertible, Lemma divisors-lemma-effective-Cartier-makes-dimension-drop; Varieties, Lemma varieties-lemma-algebraic-scheme-dim-0. Supplies the Cartier zero scheme Z of the regular section, its dimension 0, and the exact sequences 0 -> L\textasciicircum{}tensor n -> L\textasciicircum{}tensor n+1 -> i\_*(L\textasciicircum{}tensor n+1|\_Z) -> 0.
\item Proj of a section ring generated in bounded degree. Morphisms, Lemma morphisms-lemma-proper-ample-is-proj; Constructions, Lemma constructions-lemma-surjective-graded-rings-generated-degree-1-map-proj. Identifies X = Proj(S) for S the section ring and converts generation of the Veronese S\textasciicircum{}(6) in degree 1 into very-ampleness of L\textasciicircum{}tensor 6.
\end{itemize}\medskip
\ifdefined\hypertarget\hypertarget{sourcepointer-1}{}\fi
\noindent\textit{Source reference.} Original lemma 02M9: sources/stacks-chow/tags.txt, source lines 3436--3436. Prerequisite pointer invoked at proof line(s) 1538; it names a standard fact of a sibling chapter that the counted argument uses, and every step of the counted proof built on it is retained verbatim in the delivered body.\par\medskip
\ifdefined\hypertarget\hypertarget{sourcepointer-2}{}\fi
\noindent\textit{Source reference.} Original lemma 056N: sources/stacks-chow/tags.txt, source lines 6716--6716. Prerequisite pointer invoked at proof line(s) 1534; it names a standard fact of a sibling chapter that the counted argument uses, and every step of the counted proof built on it is retained verbatim in the delivered body.\par\medskip
\ifdefined\hypertarget\hypertarget{sourcepointer-3}{}\fi
\noindent\textit{Source reference.} Original section 0C4S: sources/stacks-chow/tags.txt, source lines 15237--15237. Prerequisite pointer invoked at proof line(s) 1521; it names a standard fact of a sibling chapter that the counted argument uses, and every step of the counted proof built on it is retained verbatim in the delivered body.\par\medskip
\ifdefined\hypertarget\hypertarget{sourcepointer-4}{}\fi
\noindent\textit{Source reference.} Original lemma 0C6J: sources/stacks-chow/tags.txt, source lines 15299--15299. Prerequisite pointer invoked at proof line(s) 1588; it names a standard fact of a sibling chapter that the counted argument uses, and every step of the counted proof built on it is retained verbatim in the delivered body.\par\medskip
\ifdefined\hypertarget\hypertarget{sourcepointer-5}{}\fi
\noindent\textit{Source reference.} Original lemma 06LH: sources/stacks-chow/tags.txt, source lines 8466--8466. Prerequisite pointer invoked at proof line(s) 1536; it names a standard fact of a sibling chapter that the counted argument uses, and every step of the counted proof built on it is retained verbatim in the delivered body.\par\medskip
\ifdefined\hypertarget\hypertarget{sourcepointer-6}{}\fi
\noindent\textit{Source reference.} Original lemma 0AYT: sources/stacks-chow/tags.txt, source lines 13801--13801. Prerequisite pointer invoked at proof line(s) 1541; it names a standard fact of a sibling chapter that the counted argument uses, and every step of the counted proof built on it is retained verbatim in the delivered body.\par\medskip
\begin{lemma}
\label{lemma-criterion-very-ample}
Let $k$ be a field. Let $X$ be a proper scheme over $k$ having dimension $1$
and $H^0(X, \mathcal{O}_X) = k$. Let $\mathcal{L}$ be an invertible
$\mathcal{O}_X$-module. Assume
\begin{enumerate}
\item $\mathcal{L}$ has a regular global section,
\item $H^1(X, \mathcal{L}) = 0$, and
\item $\mathcal{L}$ is ample.
\end{enumerate}
Then $\mathcal{L}^{\otimes 6}$ is very ample on $X$ over $k$.
\end{lemma}

\begin{proof}
Let $s$ be a regular global section of $\mathcal{L}$. Let
$i : Z = Z(s) \to X$ be the zero scheme of $s$, see
Divisors, Section \ref{divisors-section-effective-Cartier-invertible}.
By condition (3) we see that $Z \not = \emptyset$ (small detail omitted).
Consider the short exact sequence
$$
0 \to \mathcal{O}_X \xrightarrow{s} \mathcal{L} \to
i_*(\mathcal{L}|_Z) \to 0
$$
Tensoring with $\mathcal{L}$ we obtain
$$
0 \to \mathcal{L} \to \mathcal{L}^{\otimes 2} \to
i_*(\mathcal{L}^{\otimes 2}|_Z) \to 0
$$
Observe that $Z$ has dimension $0$
(Divisors, Lemma \ref{divisors-lemma-effective-Cartier-makes-dimension-drop})
and hence is the spectrum of an Artinian ring
(Varieties, Lemma \ref{varieties-lemma-algebraic-scheme-dim-0})
hence $\mathcal{L}|_Z \cong \mathcal{O}_Z$
(Algebra, Lemma \ref{algebra-lemma-locally-free-semi-local-free}).
The short exact sequence also shows that
$H^1(X, \mathcal{L}^{\otimes 2}) = 0$ (for example using
Varieties, Lemma \ref{varieties-lemma-chi-tensor-finite}
to see vanishing in the spot on the right). Using induction
on $n \geq 1$ and the sequence
$$
0 \to \mathcal{L}^{\otimes n} \xrightarrow{s}
\mathcal{L}^{\otimes n + 1} \to
i_*(\mathcal{L}^{\otimes n + 1}|_Z) \to 0
$$
we see that $H^1(X, \mathcal{L}^{\otimes n}) = 0$ for $n > 0$
and that there exists a global section $t_{n + 1}$
of $\mathcal{L}^{\otimes n + 1}$ which gives a trivialization of
$\mathcal{L}^{\otimes n + 1}|_Z \cong \mathcal{O}_Z$.

\medskip\noindent
Consider the multiplication map
$$
\mu_n :
H^0(X, \mathcal{L}) \otimes_k H^0(X, \mathcal{L}^{\otimes n})
\oplus
H^0(X, \mathcal{L}^{\otimes 2}) \otimes_k H^0(X, \mathcal{L}^{\otimes n - 1})
\longrightarrow
H^0(X, \mathcal{L}^{\otimes n + 1})
$$
We claim this is surjective for $n \geq 3$.
To see this we consider the short exact sequence
$$
0 \to \mathcal{L}^{\otimes n} \xrightarrow{s}
\mathcal{L}^{\otimes n + 1} \to i_*(\mathcal{L}^{\otimes n + 1}|_Z) \to 0
$$
The sections of $\mathcal{L}^{\otimes n + 1}$ coming from the left in this
sequence are in the image of $\mu_n$. On the other hand, since
$H^0(\mathcal{L}^{\otimes 2}) \to H^0(\mathcal{L}^{\otimes 2}|_Z)$
is surjective (see above) and since $t_{n - 1}$ maps to a trivialization of
$\mathcal{L}^{\otimes n - 1}|_Z$
we see that $\mu_n(H^0(X, \mathcal{L}^{\otimes 2}) \otimes t_{n - 1})$
gives a subspace
of $H^0(X, \mathcal{L}^{\otimes n + 1})$ surjecting onto the global sections of
$\mathcal{L}^{\otimes n + 1}|_Z$. This proves the claim.

\medskip\noindent
From the claim in the previous paragraph we conclude
that the graded $k$-algebra
$$
S = \bigoplus\nolimits_{n \geq 0} H^0(X, \mathcal{L}^{\otimes n})
$$
is generated in degrees $0, 1, 2, 3$ over $k$.
Recall that $X = \text{Proj}(S)$, see
Morphisms, Lemma \ref{morphisms-lemma-proper-ample-is-proj}.
Thus $S^{(6)} = \bigoplus_{n} S_{6n}$ is generated in degree $1$.
This means that $\mathcal{L}^{\otimes 6}$ is very ample as desired.
\end{proof}

\end{document}
